%==============================================================================
% PDF-2 / Section 03 — Structural Diagnosability
% Full-file replacement (RUN-P2-03)
%==============================================================================

\section{Structural diagnosability (epistemic, non-algorithmic)}
\label{sec:pdf2_structural_diagnosability}

\subsection{Problem statement: what can be known from observations}
We consider an enterprise as a continuum $K_{EA}$ with admissible internal state
space $\Omega(K_{EA})$ and boundary $\partial\Omega(K_{EA})$ (failure / phase
transition boundary). Observations do not reveal $\Omega(K_{EA})$ directly; they
are manifestations produced by interaction with an environment $E$ and an
observation operator $\Pi$.

We use the (already established) manifestation map
\[
\Sigma_{obs}(K_{EA},E) := \Pi\!\left(\Omega(K_{EA}) \cap \Omega(E)\right),
\]
and interpret any concrete dataset as a finite observation trace
$D \subseteq \Sigma_{obs}$ collected under a specific observational regime
$\Pi$ and environment $E$.

\paragraph{Diagnosability question.}
Given a finite observation trace $D$, what can be inferred about the internal
state of $K_{EA}$ (and about its proximity to $\partial\Omega(K_{EA})$) without
introducing prescriptive assumptions, control intent, or optimization criteria?

The answer is \emph{set-valued}: the data typically admits multiple internal
realizations.

\subsection{Compatibility set and non-identifiability}
\label{sec:pdf2_compatibility_set}

\begin{definition}[Compatibility set]
\label{def:compatibility_set}
Fix $(E,\Pi)$ and an observation trace $D$.
The \emph{compatibility set} is the set of all enterprise continua compatible
with the observations:
\[
\mathcal{K}_{compat}(D;E,\Pi) :=
\left\{K_{EA} \;:\; D \subseteq \Sigma_{obs}(K_{EA},E)\right\}.
\]
\end{definition}

\begin{definition}[State-level compatibility (within a fixed $K_{EA}$)]
\label{def:state_level_compatibility}
If the formalization distinguishes internal states $x \in \Omega(K_{EA})$, then
for fixed $(K_{EA},E,\Pi)$ define the set of compatible internal states:
\[
\Omega_{compat}(D \mid K_{EA},E,\Pi)
:= \left\{x \in \Omega(K_{EA}) \;:\; D \subseteq \Sigma_{obs}(x;K_{EA},E,\Pi)\right\},
\]
where $\Sigma_{obs}(x;\cdot)$ denotes the observable manifestation generated by
state $x$ under $(E,\Pi)$.
\end{definition}

\begin{definition}[Structural diagnosability (epistemic)]
\label{def:structural_diagnosability}
A diagnostic question $Q$ (e.g.\ “is the system within $\varepsilon$ of
$\partial\Omega$?” or “which diagnostic regime is admissible?”) is
\emph{structurally diagnosable} from $D$ under $(E,\Pi)$ if $Q$ is constant over
the compatibility set:
\[
\forall K_1,K_2 \in \mathcal{K}_{compat}(D;E,\Pi):\quad Q(K_1) = Q(K_2).
\]
Equivalently, for a fixed $K_{EA}$, $Q$ is diagnosable at the state level if it
is constant over $\Omega_{compat}(D \mid K_{EA},E,\Pi)$.
\end{definition}

\paragraph{Interpretation.}
Diagnosability is an epistemic property of the triple $(D,E,\Pi)$ relative to
the formal object $K_{EA}$, not a promise that an algorithm exists, nor a
guarantee of uniqueness. It is a criterion for when a \emph{definite} diagnostic
statement is justified.

\subsection{Underdetermination as the default}
The manifestation map is generally non-injective:
different internal continua (or different internal states) can produce
indistinguishable observation traces under the same $(E,\Pi)$.

\begin{proposition}[Generic non-identifiability]
\label{prop:generic_nonidentifiability}
Unless $\Pi$ is informationally complete for $\Omega(K_{EA})$ (a degenerate,
rarely realizable condition in enterprise contexts), the mapping
$K_{EA} \mapsto \Sigma_{obs}(K_{EA},E)$ is not injective. Therefore,
$\lvert\mathcal{K}_{compat}(D;E,\Pi)\rvert > 1$ is generic for finite $D$.
\end{proposition}

\begin{proof}[Proof sketch]
Observation operators $\Pi$ reduce dimension: $\dim \Sigma_{obs} \ll \dim
\Omega(K_{EA})$ in realistic settings. Dimension reduction together with finite
sampling implies that multiple distinct internal structures can share the same
projection. Hence, non-identifiability is generic.
\end{proof}

\paragraph{Negative result (accepted).}
It is generally impossible to \emph{reconstruct} the internal structure of an
enterprise continuum from observational traces. The maximal justified output of
diagnosis is a \emph{set} (compatibility set) and statements invariant over that
set.

\subsection{Diagnosable properties: invariants and bounds}
Even under non-identifiability, certain \emph{invariants} and \emph{bounds} can
be diagnosable.

\begin{definition}[Diagnosable invariant]
\label{def:diagnosable_invariant}
A functional $I$ defined on continua is a \emph{diagnosable invariant} from $D$
if it is constant on $\mathcal{K}_{compat}(D;E,\Pi)$.
\end{definition}

\begin{definition}[Diagnosable bound]
\label{def:diagnosable_bound}
A scalar functional $B$ is \emph{diagnosably bounded} from $D$ if there exist
functions $b^-(D)$ and $b^+(D)$ such that for all
$K \in \mathcal{K}_{compat}(D;E,\Pi)$,
\[
b^-(D) \le B(K) \le b^+(D),
\]
and the bounds depend only on $D$ (and fixed $(E,\Pi)$).
\end{definition}

\paragraph{Typical examples (structural, not managerial).}
Examples of candidates for $B$ include:
distance-to-boundary proxies, admissibility of certain transitions $\Psi$, or
whether the feasible model space $M_{feasible}$ is empty. These are not
performance targets; they are \emph{structural feasibility} statements.

\subsection{Model-space triad and the STOP boundary}
We assume the established triad of model spaces:
$M_{current}$ (what is), $M_{declared}$ (what is claimed), and $M_{feasible}$
(what can exist without violating thresholds / constraints). A central negative
outcome is the collapse of feasibility.

\begin{definition}[Feasibility collapse (STOP condition)]
\label{def:feasibility_collapse_stop}
A trace $D$ implies a \emph{feasibility collapse} (diagnostic STOP) under
$(E,\Pi)$ if
\[
\forall K \in \mathcal{K}_{compat}(D;E,\Pi):\quad M_{feasible}(K) = \varnothing.
\]
\end{definition}

\paragraph{Meaning.}
STOP is not advice. It is a formal diagnostic classification: within the
compatibility set supported by $D$, there exists no internally consistent model
meeting the declared constraints. It is an epistemic statement about
non-existence.

\subsection{Diagnosability of proximity to $\partial\Omega$}
The boundary $\partial\Omega(K_{EA})$ is generally not observable directly.
However, one can ask whether closeness to the boundary is diagnosable as an
invariant over compatible continua/states.

\begin{definition}[Diagnosable proximity statement]
\label{def:diagnosable_proximity}
Let $d(\cdot,\partial\Omega)$ be a distance-like functional defined on internal
states $x \in \Omega(K_{EA})$ (or on continua via induced metrics, if defined in
the core). A statement of the form
\[
d(\cdot,\partial\Omega) \le \varepsilon
\]
is diagnosable from $D$ if it holds for all compatible realizations (continua or
states, depending on the level of formalization).
\end{definition}

\paragraph{Negative result (accepted).}
Without strong assumptions on $\Pi$ and on the metric structure, proximity to
$\partial\Omega$ is usually \emph{not} diagnosable; at best, one obtains
one-sided bounds.

\subsection{Epistemic closure: what this section does \emph{not} claim}
This section does not claim:
(i) unique identification of internal structure;
(ii) existence of a computable estimator;
(iii) controllability, optimizability, or policy implications.

It establishes only the epistemic formalism required to speak precisely about
what is (and is not) inferable from observations in the $K_{EA}$ setting.
